\chapter{Kiểm Thử và Đánh Giá}

\section{Chiến lược kiểm thử}

Hệ thống được kiểm thử theo ba cấp độ:

\begin{description}
  \item[Unit Testing:] Kiểm thử các hàm độc lập (Haversine, NMEA parser, alert state machine) bằng \texttt{pytest}.
  \item[Integration Testing:] Kiểm thử luồng dữ liệu end-to-end từ MQTT $\rightarrow$ Backend $\rightarrow$ MongoDB $\rightarrow$ WebSocket.
  \item[System Testing:] Kiểm thử toàn bộ hệ thống với thiết bị thực tế ngoài trời.
\end{description}

\section{Kiểm thử đơn vị}

\subsection{Kiểm thử thuật toán Haversine}

\begin{lstlisting}[language=python, caption=Unit test thuật toán Haversine]
def test_haversine_known_distance():
    # Hà Nội (Hoàn Kiếm) <-> Hà Nội (Cầu Giấy): ~5.2 km
    d = haversine(21.0285, 105.8542, 21.0367, 105.7821)
    assert abs(d - 5200) < 100   # sai số < 100m

def test_haversine_same_point():
    d = haversine(21.0285, 105.8542, 21.0285, 105.8542)
    assert d == 0.0

def test_haversine_geofence_boundary():
    # Điểm tại chính xác 100m khỏi tâm
    lat_offset = 100 / 111320  # ~0.0009 độ
    d = haversine(21.0285, 105.8542, 21.0285 + lat_offset, 105.8542)
    assert abs(d - 100) < 1     # sai số < 1m
\end{lstlisting}

\subsection{Kiểm thử máy trạng thái cảnh báo}

\begin{table}[H]
\centering
\caption{Kết quả kiểm thử các chuyển trạng thái cảnh báo}
\label{tab:alert_state_test}
\begin{tabularx}{\textwidth}{l l l X l}
\toprule
\textbf{TC} & \textbf{Trạng thái cũ} & \textbf{GPS mới} & \textbf{Điều kiện} & \textbf{Kết quả} \\
\midrule
TC-01 & inside  & outside & --              & \texttt{outside\_entered} \\
TC-02 & outside & inside  & --              & \texttt{inside\_entered} \\
TC-03 & outside & outside & $t <$ cooldown  & \texttt{None} \\
TC-04 & outside & outside & $t \geq$ cooldown & \texttt{outside\_reminder} \\
TC-05 & inside  & inside  & --              & \texttt{None} \\
TC-06 & inside  & outside & nhiễu $<$5\,m   & \texttt{None} (lọc) \\
\bottomrule
\end{tabularx}
\end{table}

Tất cả 6 test case đều pass, xác nhận logic máy trạng thái hoạt động đúng đặc tả.

\section{Kiểm thử tích hợp}

\subsection{Kiểm thử pipeline MQTT -- Backend -- WebSocket}

Sử dụng script Python (\texttt{test\_mqtt\_client.py}) để giả lập thiết bị GPS gửi dữ liệu:

\begin{lstlisting}[language=python, caption=Script giả lập thiết bị GPS]
import paho.mqtt.client as mqtt, json, time, random

def publish_test_locations():
    client = mqtt.Client()
    client.connect("localhost", 1883)
    for i in range(100):
        payload = {
            "deviceId": "child_01",
            "lat": 21.0285 + random.uniform(-0.001, 0.001),
            "lng": 105.8542 + random.uniform(-0.001, 0.001),
            "timestamp": int(time.time())
        }
        client.publish("gps/child_01", json.dumps(payload))
        time.sleep(1)
\end{lstlisting}

\textbf{Kết quả:} 100/100 tin nhắn được lưu vào MongoDB, WebSocket broadcast thành công đến 3 client Flutter đang kết nối đồng thời. Không có tin nhắn bị mất.

\subsection{Kiểm thử API REST}

\begin{table}[H]
\centering
\caption{Kết quả kiểm thử các endpoint REST}
\label{tab:api_test}
\begin{tabularx}{\textwidth}{l l l X}
\toprule
\textbf{Endpoint} & \textbf{Method} & \textbf{Status} & \textbf{Thời gian phản hồi} \\
\midrule
\texttt{/auth/login}          & POST   & 200 & 120\,ms \\
\texttt{/latest/child\_01}    & GET    & 200 & 45\,ms \\
\texttt{/history/child\_01}   & GET    & 200 & 180\,ms (100 điểm) \\
\texttt{/geofence/update}     & POST   & 200 & 55\,ms \\
\texttt{/stats/child\_01}     & GET    & 200 & 320\,ms \\
\texttt{/stats/child\_01/heatmap} & GET & 200 & 450\,ms \\
\bottomrule
\end{tabularx}
\end{table}

Tất cả endpoint hoạt động đúng, thời gian phản hồi nằm trong ngưỡng yêu cầu ($<$500\,ms).

\section{Kiểm thử hệ thống thực tế}

\subsection{Điều kiện kiểm thử}

\begin{itemize}[noitemsep]
  \item \textbf{Địa điểm:} Khuôn viên Đại học Bách khoa Hà Nội (khu ngoài trời).
  \item \textbf{Thời gian:} Buổi sáng, trời quang, góc mặt trời cao.
  \item \textbf{Thiết bị:} ESP32 + SIM7000G (SIM Viettel 4G), pin 2000\,mAh.
  \item \textbf{Backend:} Chạy trên Render.com, MongoDB Atlas.
  \item \textbf{App:} Flutter Android trên điện thoại Samsung Galaxy S23.
\end{itemize}

\subsection{Độ chính xác GPS}

Để đánh giá độ chính xác, thiết bị được đặt cố định tại một vị trí biết tọa độ chính xác trong 10 phút (120 điểm GPS), rồi tính sai số:

\begin{table}[H]
\centering
\caption{Kết quả đo độ chính xác GPS (120 mẫu, ngoài trời)}
\label{tab:gps_accuracy}
\begin{tabular}{l r}
\toprule
\textbf{Chỉ số} & \textbf{Giá trị} \\
\midrule
Sai số trung bình (CEP 50\%) & 3.2\,m \\
Sai số cực đại (CEP 95\%) & 6.8\,m \\
Độ lệch chuẩn           & 1.4\,m \\
Tỷ lệ điểm sai $<$5\,m  & 87\% \\
Tỷ lệ điểm sai $<$10\,m & 99\% \\
\bottomrule
\end{tabular}
\end{table}

\textbf{Nhận xét:} Độ chính xác đạt 3.2\,m CEP trong điều kiện ngoài trời thoáng đáp ứng yêu cầu đề ra ($<$5\,m). Độ lọc nhiễu 5\,m loại bỏ hiệu quả dao động GPS tự nhiên của thiết bị đứng yên.

\subsection{Độ trễ end-to-end}

Đo thời gian từ khi thiết bị gửi MQTT đến khi Flutter app cập nhật vị trí trên bản đồ:

\begin{table}[H]
\centering
\caption{Đo độ trễ end-to-end (50 lần đo)}
\label{tab:latency}
\begin{tabular}{l r}
\toprule
\textbf{Chỉ số} & \textbf{Giá trị} \\
\midrule
Độ trễ trung bình & 1.8\,giây \\
Độ trễ p50 (median) & 1.5\,giây \\
Độ trễ p95 & 3.2\,giây \\
Độ trễ tối đa & 4.7\,giây \\
\bottomrule
\end{tabular}
\end{table}

\textbf{Phân tích:} Tổng độ trễ bao gồm: thời gian gửi MQTT qua GPRS ($\sim$0.5\,s) + xử lý backend ($\sim$0.1\,s) + truyền WebSocket ($\sim$0.1\,s) + render Flutter ($\sim$0.05\,s). Tổng $<$1\,s trong điều kiện mạng tốt. Trường hợp p95 tăng lên 3.2\,s do biến động mạng GPRS.

\subsection{Kiểm thử cảnh báo geofence}

\begin{table}[H]
\centering
\caption{Kết quả kiểm thử cảnh báo geofence (bán kính 50\,m)}
\label{tab:geofence_test}
\begin{tabularx}{\textwidth}{l X r r}
\toprule
\textbf{Kịch bản} & \textbf{Mô tả} & \textbf{Kết quả} & \textbf{Thời gian} \\
\midrule
KT-01 & Rời khỏi vùng -- email & Gửi thành công & 3.1\,s \\
KT-02 & Rời khỏi vùng -- Telegram & Gửi thành công & 2.7\,s \\
KT-03 & Trở về vùng an toàn & Thông báo đúng & 2.9\,s \\
KT-04 & Dao động ở biên giới & Không gửi (cooldown) & -- \\
KT-05 & Vùng di động theo phụ huynh & Cập nhật đúng & $<$2\,s \\
KT-06 & Đa thiết bị (3 thiết bị) & Tất cả nhận cảnh báo & 3.5\,s \\
\bottomrule
\end{tabularx}
\end{table}

\subsection{Kiểm thử thời lượng pin}

Thiết bị chạy liên tục với chu kỳ gửi GPS 5 giây/lần, pin LiPo 2000\,mAh:

\begin{table}[H]
\centering
\caption{Đo thời lượng pin}
\label{tab:battery}
\begin{tabular}{l r}
\toprule
\textbf{Chế độ} & \textbf{Thời lượng} \\
\midrule
Hoạt động đầy đủ (GPS + GPRS + Publish 5s) & 9.5 giờ \\
Giảm tần suất publish (30 giây) & $\sim$24 giờ \\
Deep sleep giữa các lần publish & $>$48 giờ (ước tính) \\
\bottomrule
\end{tabular}
\end{table}

\section{Tổng kết đánh giá}

\begin{table}[H]
\centering
\caption{So sánh kết quả đạt được với yêu cầu đề ra}
\label{tab:evaluation_summary}
\begin{tabularx}{\textwidth}{l X r r}
\toprule
\textbf{Chỉ số} & \textbf{Yêu cầu} & \textbf{Đạt được} & \textbf{Kết quả} \\
\midrule
Độ chính xác GPS   & $<$5\,m (CEP 50\%)  & 3.2\,m  & \checkmark \\
Độ trễ end-to-end  & $<$5\,s             & 1.8\,s  & \checkmark \\
Thời gian phản hồi API & $<$500\,ms      & $<$450\,ms & \checkmark \\
Thời lượng pin     & $\geq$8 giờ         & 9.5 giờ & \checkmark \\
Cảnh báo email     & $<$5\,s             & 3.1\,s  & \checkmark \\
Hỗ trợ đa thiết bị & $\geq$3             & Kiểm thử 3 & \checkmark \\
\bottomrule
\end{tabularx}
\end{table}

Tất cả các chỉ số đánh giá đều đạt hoặc vượt yêu cầu đề ra ban đầu. Hệ thống hoạt động ổn định trong điều kiện kiểm thử thực tế.
